% Mayor Climate Game - APA 7 student manuscript
\documentclass[stu, 12pt, apacite]{apa7}

\title{Mayor Climate Game: A Browser-Based Prototype for Teaching Urban Energy and Ventilation Tradeoffs}
\shorttitle{Mayor Climate Game}

\author{Hu Yunhan}
\affiliation{Generation AI 2026 Research Project}

\abstract{
Urban planners and citizens face coupled decisions about energy cost, carbon emissions, and ventilation-sensitive siting, yet professional simulation tools remain inaccessible to many learners. This student paper reports a pilot evaluation of \textit{Mayor Climate Game}, a browser-based serious game in which participants allocate a fixed municipal budget across coal, wind, and ground-source heat pump facilities on a terrain-typed grid while monitoring live climate indicators. The research question asks how the artifact shapes facility choices and post-play understanding of budget--carbon--ventilation tradeoffs among undergraduate participants. We conducted an exploratory pilot in July 2026 with three completed sessions logged to Supabase (two distinct facility layouts). All three participants answered the ventilation knowledge item correctly and reported choosing coal to save budget. Observed final scores were 84 and 70 (composite client-side metric), with total carbon emissions of 120 and 600\,t CO\textsubscript{2}e, respectively; no session placed coal on wind outlets. Findings are interpreted cautiously as preliminary patterns rather than confirmatory evidence.
}

\begin{document}
\maketitle

\section{Introduction}

Cities confront rising heat, cooling-related energy demand, and pressure to reduce greenhouse-gas emissions \cite{iea2023stayingcool,iea2023sustainablecooling}. Urban ventilation corridors and wind outlets can help dissipate heat and pollution when land-use and energy siting respect airflow pathways \cite{zheng2022urban,guan2015guiyang}. Professional models (e.g., CFD, WRF, GIS workflows) capture these dynamics with fidelity but require expertise and infrastructure beyond typical classroom settings. Serious games offer a complementary route: learners manipulate simplified systems, receive immediate feedback, and reveal preferences through logged choices \cite{seriousgame2019neighbourhood}.

\textit{Mayor Climate Game} was developed as a research prototype that combines rule-based microclimate feedback with automatic session recording. The present report documents an early pilot after deployment to a public URL.

\textbf{Research question.} How does \textit{Mayor Climate Game} affect energy layout choices and post-play understanding of tradeoffs among cost, carbon emissions, and urban ventilation for undergraduate participants in a social-science or urban-climate learning context?

\textbf{Hypothesis.} After playing one round with live metrics and advisory cues, participants will place fewer high-carbon coal facilities in ventilation-sensitive cells and will show accurate understanding of the budget--carbon--ventilation tradeoff on a brief post-play survey, relative to what would be expected if cost alone drove choices.

\section{Method}

\subsection{Artifact}

The artifact is a single-page web application (Vercel-hosted; local \texttt{index.html} backup) requiring no installation. Participants act as a mayor with CNY\,10 million, placing coal power (lower capital cost, higher carbon), wind power, or ground-source heat pumps on a 12\,$\times$\,12 grid with normal zones, residential areas, ventilation corridors, wind outlets, and low basins.

The current interface includes a \textbf{mission bar} that states session goals and step-level guidance (e.g., prompting facility selection), a \textbf{live score} panel alongside real-time average temperature, wind speed, and total carbon emissions, and post-session export of a human-readable \textbf{TXT report} (layout, metrics, score, and survey responses). When server credentials are configured, download triggers a client POST to \texttt{/api/submit}, which stores session rows in \textbf{Supabase} (\texttt{public.game\_results}: budget used, score, metrics, layout JSON, survey JSON, and raw payload). Rule-based advisor text (not a live large language model during gameplay) flags risky placements such as coal at wind outlets.

The authors used an AI assistant to convert the manuscript from Markdown into APA 7 \LaTeX{} and to help fix compilation errors. All research content, preliminary data interpretation, and references were produced and verified by the author.

\subsection{Material Boundaries}

In-game temperature, wind, and carbon values are heuristic functions of facility parameters, terrain modifiers, and placement locations---not outputs of CFD, WRF, or GIS simulation. Step-level impact cards and advisor messages are deterministic scripts. An optional post-play ``generate AI summary'' button, when enabled on the server, calls an external chat API after gameplay only; it did not drive in-game mechanics in the pilot build.

\subsection{Design}

Each session comprised one unrestricted placement phase, continuous metric feedback, and two required closed-ended survey items administered in a finish modal: (1) whether coal was chosen to save budget (\texttt{budgetChoice}: yes/no) and (2) whether coal at a wind outlet reduces city-wide ventilation (\texttt{windUnderstanding}: correct/incorrect). Behavioral variables (final \texttt{score}, \texttt{metrics.carbon}, facility counts, terrain-specific placements) were derived from exported layout JSON and Supabase records.

\subsection{Participants}

Three volunteers recruited in July 2026 through the Generation AI 2026 Research Project course context completed one round each and submitted data captured in Supabase (three sessions total). Participation was voluntary; identifiers were not collected in the stored gameplay payload beyond timestamps and layout records.

\subsection{Coding}

Survey items were coded dichotomously (yes/correct = 1). Layout arrays were parsed programmatically to count facility types and to flag coal on \texttt{outlet} or \texttt{corridor} cells. Sessions with identical facility coordinates were treated as one unique layout for descriptive reporting.

\subsection{Limitations}

The microclimate model is simplified; inferential claims are exploratory. The pilot sample ($N = 3$) lacks power for generalization. Self-reported budget motivation may not align with observed placement when multiple facilities are used. Network or Supabase configuration failures could omit server-side rows even when local TXT download succeeds---all three pilot sessions reported here were verified in Supabase.

\section{Results}

Three sessions were recorded in July 2026. Post-play survey performance was uniform: \textbf{100\%} of participants answered the ventilation item correctly (3/3) and \textbf{100\%} reported \texttt{budgetChoice} = yes (3/3), indicating that all attributed coal use partly to cost savings while still endorsing the ventilation rule verbally.

Two unique facility layouts emerged across the three sessions (one layout appeared twice). For \textbf{layout 1} ($n = 2$ sessions), the final composite \textbf{score} was \textbf{84} with total carbon emissions of \textbf{120\,t} CO\textsubscript{2}e. For \textbf{layout 2} ($n = 1$ session), the final \textbf{score} was \textbf{70} with total carbon emissions of \textbf{600\,t} CO\textsubscript{2}e. Parsing placement coordinates showed \textbf{no coal facilities on wind outlets} in any session (0/3).

These descriptive contrasts suggest that lower-carbon layouts co-occurred with higher composite scores in this pilot, while the highest carbon layout still avoided outlet siting that would trigger strong in-game ventilation penalties. No inferential tests were conducted given $N = 3$.

\section{Discussion}

Pilot data align partially with the hypothesis: participants demonstrated accurate ventilation knowledge and avoided coal on outlets, yet uniformly reported budget-driven coal motivation---consistent with coal's lower capital cost in the instrument. The spread in carbon outcomes (120 vs.\ 600\,t) between layouts indicates that cost-focused strategies can still diverge sharply in emissions depending on facility mix, supporting the game's design goal of making tradeoffs visible \cite{iea2023stayingcool}.

Limitations of model fidelity and sample size require reframing these patterns as hypotheses for a larger study. Future work should pre-register primary outcomes, expand $N$, and compare rule-based advisory conditions. Connecting logged layouts to ventilation literature \cite{zheng2022urban,guan2015guiyang} and serious-game learning mechanisms \cite{seriousgame2019neighbourhood} will clarify whether short browser sessions shift siting behavior or mainly assess declarative understanding, as this exploratory pilot begins to suggest.

\bibliography{references}

\end{document}
